\section{Ciclo cellulare}

Il ciclo cellulare si articola in:
\begin{enumerate}
  \item \textbf{interfase}: fase G$_1$, fase S, fase G$_2$;
  \item \textbf{fase M}: mitosi (profase, prometafase, metafase, anafase, telofase) e citocinesi.
\end{enumerate}

% ---------------------------------------------------------------------------
\subsection{Interfase}

Periodo che va dalla conclusione di una mitosi all'inizio della successiva.
\begin{itemize}
  \item \textbf{fase G$_1$}: sintesi di proteine e altre molecole, ma non di DNA; la cellula cresce e possiede un normale metabolismo, gli organelli si duplicano;
  \item \textbf{fase S}: sintesi del DNA e delle proteine cromosomiche $\rightarrow$ si replica il DNA e si duplicano i cromosomi;
  \item \textbf{fase G$_2$}: la cellula cresce e si prepara alla mitosi.
\end{itemize}
Una cellula può uscire dal ciclo ed entrare in \textbf{fase G$_0$}: arresto del ciclo cellulare.

% ---------------------------------------------------------------------------
\subsection{Mitosi}

\rubrica{Profase}
\begin{enumerate}
  \item il materiale cromosomico si condensa per formare cromosomi mitotici compatti, costituiti da due cromatidi uniti fra loro a livello del centromero;
  \item il citoscheletro scompare e si assembla il fuso mitotico;
  \item il complesso di Golgi e il reticolo endoplasmatico si frammentano; l'involucro nucleare si disperde.
\end{enumerate}

\rubrica{Prometafase}
\begin{enumerate}
  \item i microtubuli cromosomici si collegano al cinetocore dei cromosomi;
  \item i cromosomi si muovono verso l'equatore del fuso.
\end{enumerate}

\rubrica{Metafase}
I cromosomi si allineano sulla piastra metafasica, connessi ad entrambi i poli tramite i microtubuli cromosomici.

\rubrica{Anafase}
\begin{enumerate}
  \item i centromeri si dividono e i cromatidi si separano;
  \item i cromosomi si muovono verso i poli opposti del fuso;
  \item i poli del fuso si allontanano in direzione opposta.
\end{enumerate}

\rubrica{Telofase}
\begin{enumerate}
  \item i cromosomi si raggruppano ai poli opposti;
  \item i cromosomi si disperdono;
  \item l'involucro nucleare si riassembla intorno ai cromosomi raggruppati;
  \item si riformano il complesso di Golgi e il reticolo endoplasmatico;
  \item le cellule si separano per mezzo della citocinesi.
\end{enumerate}

% ---------------------------------------------------------------------------
\subsection{Citocinesi}

\rubrica{Cellule animali}
Si forma un \textbf{anello contrattile di microfilamenti}: il solco inizia come una dentellatura che corre lungo tutto il perimetro della cellula in corrispondenza del piano della regione mediana del precedente fuso; il solco si fa più profondo in seguito alla contrazione dei microfilamenti, come un elastico che si stringe intorno alla cellula, finché i due nuclei figli non sono separati in cellule distinte.

\rubrica{Cellule vegetali}
Uno strato di vescicole contenenti materiale della parete cellulare si raccoglie sul piano della regione mediana del fuso precedente; sempre più vescicole si aggiungono finché lo strato si estende attraverso tutta la cellula; le vescicole si fondono insieme, liberando il loro contenuto all'interno di una parete che gradualmente si accresce tra le due cellule figlie, la \textbf{piastra cellulare}, finché le cellule figlie non sono separate da una nuova parete continua.

% ---------------------------------------------------------------------------
\subsection{Il fuso mitotico}

\begin{itemize}
  \item durante l'interfase, il centrosoma contiene una coppia di centrioli;
  \item la coppia originale di centrioli si duplica;
  \item il centrosoma si divide in due parti, ognuna contenente un centriolo parentale e un centriolo figlio;
  \item i centrosomi si muovono ai lati opposti del nucleo, mantenendosi collegati mediante una gran quantità di microtubuli, che formano il fuso precoce.
\end{itemize}

\rubrica{Componenti del fuso completo}
\begin{itemize}
  \item \textbf{microtubuli del cinetocore} (microtubuli del fuso): collegano i cinetocori ai poli;
  \item \textbf{microtubuli astrali}: si irradiano in ogni direzione formando l'\textbf{aster};
  \item \textbf{microtubuli polari}: si sovrappongono sul piano equatoriale;
  \item \textbf{piastra metafasica}: piano equatoriale della cellula.
\end{itemize}

% ---------------------------------------------------------------------------
\subsection{Meiosi}

Si verifica soltanto nelle cellule germinali degli animali e delle piante a riproduzione sessuata. Consiste in 2 divisioni consecutive (meiosi I e meiosi II) accompagnate da una sola duplicazione del DNA. Il risultato è la formazione di 4 cellule aploidi, con corredo cromosomico dimezzato $\rightarrow$ i \textbf{gameti}.

\rubrica{Interfase}
Corrisponde all'interfase delle cellule somatiche.
\begin{itemize}
  \item ha una durata di 20--30 ore;
  \item si ha la duplicazione del DNA e il raddoppiamento di ciascun cromosoma in 2 cromatidi.
\end{itemize}

% ---------------------------------------------------------------------------
\subsection{Meiosi I}

\rubrica{Profase I}
\begin{enumerate}
  \item i cromosomi omologhi si appaiano punto per punto (\textit{sinapsi});
  \item i cromosomi iniziano a condensarsi; appaiono evidenti le coppie di cromosomi omologhi, e ogni omologo è formato da 2 cromatidi fratelli adiacenti $\rightarrow$ \textbf{tetradi} (4 cromatidi); si verifica il \textbf{crossing-over}.
\end{enumerate}

\rubrica{Metafase I}
Le tetradi si allineano sul piano equatoriale della cellula e restano unite a livello dei \textbf{chiasmi} (siti in cui è avvenuto il crossing-over).

\rubrica{Anafase I}
I cromosomi omologhi si separano e migrano ai poli opposti; i cromatidi fratelli restano uniti a livello dei centromeri.

\rubrica{Telofase I}
Un solo cromosoma di ogni coppia di omologhi raggiunge ciascun polo. Avviene la citocinesi.

% ---------------------------------------------------------------------------
\subsection{Meiosi II}

\rubrica{Profase II}
I cromosomi si condensano nuovamente. Il DNA non si replica.

\rubrica{Metafase II}
I cromosomi si allineano lungo il piano equatoriale della cellula.

\rubrica{Anafase II}
I cromatidi fratelli si separano e migrano ai poli opposti.

\rubrica{Telofase II}
Si formano i nuclei ai poli opposti di ciascuna cellula. Avviene la citocinesi. Vengono così prodotti 4 gameti (negli animali) o 4 spore (nelle piante).

% ---------------------------------------------------------------------------
\subsection{Gametogenesi}

\rubrica{Spermatogenesi}
Nel maschio, la meiosi avviene prima del differenziamento.
\begin{itemize}
  \item gli spermatogoni proliferano per mitosi;
  \item crescita $\rightarrow$ spermatocita primario;
  \item divisioni meiotiche: spermatocita primario $\rightarrow$ 2 spermatociti secondari $\rightarrow$ 4 spermatidi;
  \item differenziazione: 4 spermatidi $\rightarrow$ 4 spermatozoi.
\end{itemize}

\rubrica{Oogenesi}
Nella femmina, la meiosi avviene dopo il differenziamento.
\begin{itemize}
  \item gli oogoni proliferano per mitosi;
  \item crescita e differenziamento $\rightarrow$ oocita primario;
  \item meiosi: oocita primario $\rightarrow$ oocita secondario $+$ primo corpo polare;
  \item fertilizzazione: oocita secondario $\rightarrow$ cellula uovo $+$ secondo corpo polare;
  \item ogni oocita primario forma quindi un singolo uovo fecondabile e due o tre corpuscoli polari.
\end{itemize}

% ---------------------------------------------------------------------------
\subsection{Mitosi e meiosi a confronto}

Confronto tra gli eventi delle due divisioni, a partire da una cellula diploide con 4 cromosomi (2 paia di omologhi):

\begin{table}[htbp]
  \centering
  \small
  \renewcommand{\arraystretch}{1.3}
  \begin{tabularx}{\textwidth}{|l|X|X|}
    \hline
    \textbf{Tappa} & \textbf{Mitosi} & \textbf{Meiosi} \\
    \hline
    Profase & nessuna sinapsi dei cromosomi omologhi & sinapsi dei cromosomi omologhi per formare le tetradi (profase I) \\
    \hline
    Anafase & i cromatidi fratelli migrano ai poli opposti & i cromosomi omologhi migrano ai poli opposti (anafase I) \\
    \hline
    Profase II & -- & i cromosomi si condensano di nuovo, con cromosomi dicromatidici \\
    \hline
    Anafase II & -- & i cromatidi fratelli migrano ai poli opposti \\
    \hline
    Cellule risultanti & due cellule $2n$ con cromosomi monocromatidici & quattro cellule $n$ con cromosomi monocromatidici \\
    \hline
  \end{tabularx}
\end{table}
